\chapter{Chemistry Hides in Kitchens, Medicine Cabinets, and Gardens}

\begin{kaichang}
There is no need to visit a laboratory. Today we will do just three things, all in the kitchen.

First: find a transparent glass, fill it halfway with warm water, add a spoonful of white sugar, and stir. The sugar disappears---the water is still clear and transparent, but a sip tastes sweet.

Second: put a spoonful of white sugar in a clean small saucepan and heat it slowly over low heat (ask a parent to do this step, or watch a video; melted sugar is extremely hot and can cause serious burns if it touches skin). After a few minutes, the white sugar first melts into a clear liquid, then turns golden, then dark brown, releasing a toasted aroma. Taste a little caramel after it has cooled: fragrant and bitter, it tastes nothing like white sugar.

Third: sprinkle a small pinch of dried baking yeast into a glass of warm sugar water, stir gently, and set it aside. Come back half an hour later: a layer of foam floats on the surface, and tiny bubbles cover the sides, like water slowly coming to a boil---except it is cool. Smell it up close, and you will notice a faint alcoholic smell.

The same spoonful of sugar, three different experiences. Make a guess: in which cases has the sugar merely ``changed its appearance,'' and in which has it really ``become something else''? And how does yeast make sugar water bubble? It is alive; what does something done by a living thing have to do with chemistry? Write down your guesses without peeking at the answers ahead. At the end of this chapter, we will return with a completely new pair of glasses to examine them.
\end{kaichang}

\section{Start with Only What We Can Observe}

Science begins with honest observation, not terminology. Before introducing any concepts, let us faithfully record everything we can see, smell, taste, or measure in these three events, without a word of explanation.

\textbf{Sugar water.} When sugar enters water, its grains become smaller and smaller until they disappear completely. The water remains transparent, its color does not change, and there are no bubbles. Touch the glass: its temperature has barely changed. Yet the water is clearly sweet. Where did the sugar go? If you are patient enough to leave the sugar water somewhere well ventilated while the water slowly evaporates, after several days a white solid will reappear at the bottom. Taste it---still sugar. The sugar seems merely to have ``hidden'' in the water, ready to be recovered.

\textbf{Caramel.} Heated sugar tells a different story. Its color changes from white to gold to brownish black, and its smell from nothing to a toasted aroma. Break apart the cooled solid, and its texture, taste, and smell have all changed. More troubling still: some of the caramel dissolves in water, but no amount of evaporation or cooling brings back the original white sugar. Water vapor also escapes above the saucepan during heating, and cooking it too long may even produce choking smoke. The original sugar really seems to have disappeared, replaced by a cast of new characters.

\textbf{Sugar water with yeast.} The most conspicuous sign here is bubbling: gas is constantly being produced in the glass. Where does it come from? Leave an identical glass of sugar water without yeast beside it for half a day, and nothing happens---so the bubbles are neither produced by the water alone nor something that was hiding in the glass beforehand. Look at the other signs: the glass feels slightly warm; the sweetness fades day by day; after a day or two, a faint alcoholic taste appears. Sugar is decreasing, something new is increasing, and an invisible, intangible ``worker''---yeast---is working overtime.

That is the end of our observations. Notice that none of these three descriptions used a chemistry term, yet you probably already feel that the events are ``different in kind'': the first seems reversible, while the other two do not. Treasure that intuition; it is valuable. The task of this chapter is to find a defensible criterion for deciding whether things are ``the same or different.''

\section{Zooming In to the Particles}

Staring at the glass will never explain why a change can or cannot be reversed. We need to zoom in, to a scale ten thousand times smaller than the width of a hair, then another ten thousand times smaller still.

Beginning in Chapter 6, this book will tell the full story of atoms: how evidence gradually forced people to accept them. For now, provisionally accept the following picture as a map borrowed on credit, to be paid for later:

\begin{itemize}[itemsep=2pt]
\item All matter in the world consists of extraordinarily tiny particles;
\item Substances with their own distinctive ``personalities,'' such as sugar, water, alcohol, and carbon dioxide, each correspond to a particular kind of \textbf{molecule}---the smallest unit that retains the substance's properties;
\item Molecules are small structures made by joining \textbf{atoms} in particular ways. The selection of building blocks is surprisingly small: sugar molecules contain only three kinds of atoms, carbon, hydrogen, and oxygen, but the particular arrangement makes them sugar rather than something else.
\end{itemize}

Use this picture to revisit our three events.

\textbf{Sugar dissolving in water.} The sugar molecules themselves do not change. They simply separate from the tightly packed grains and spread evenly into the spaces among water molecules, like a box of building blocks taken apart and scattered into a large basin of marbles---not a single block is damaged. That is why the sweetness remains, and why evaporating the water lets the sugar molecules gather again and become white sugar once more.

\textbf{Sugar caramelizing when heated.} At high temperatures, the internal arrangements of sugar molecules break apart. Some atoms shed the components of ``water'' and link together into larger, darker molecules (which produce the toasted aroma and brownish black color), while others drift away as water vapor and small molecules. This time the blocks have not only been taken apart but rebuilt into completely different shapes. Can they reassemble themselves into ``sugar''? There is no simple way.

\textbf{Yeast fermentation.} Yeast is a microorganism---a real, living cell. It ``swallows'' sugar molecules, breaks them apart inside the cell, and reassembles them into two kinds of things: carbon dioxide molecules (the source of the bubbles) and alcohol molecules (the source of the alcoholic smell). Breaking down sugar releases energy that supports the yeast's growth and reproduction, incidentally warming the glass a little. Notice one detail: yeast does not create the carbon dioxide in the bubbles ``out of nothing.'' Its carbon and oxygen atoms were originally locked inside sugar molecules. Yeast takes things apart and reassembles them; it does not conjure atoms from thin air.

Our three events now fall into two clear groups: in the first, the molecules themselves stay the same; in the other two, molecules are broken apart and reorganized, producing new molecules that were not present before. This dividing line is the most fundamental one in chemistry.

\section{Physical and Chemical Changes}

We can now name the two groups.

\begin{itemize}[itemsep=2pt]
\item \textbf{Physical change}: a substance changes its shape, state, or position, but its molecules remain unchanged. Sugar dissolving in water, water freezing, alcohol evaporating, an iron wire bent into a circle---the particles are still the same particles.
\item \textbf{Chemical change} (also called a chemical reaction): old molecules are broken apart, and atoms recombine to form new molecules. Sugar caramelizing or fermenting, iron rusting, wood burning, milk souring---after the change, there are kinds of molecules that were not present before.
\end{itemize}

There is only one true criterion for classifying a change: \textbf{has a new substance (new molecules) formed?} Whether the change is violent, gives off light or heat, produces bubbles, or can be reversed is only supporting evidence, never the criterion itself.

Use this criterion to take a quick inventory of an ordinary day. When you boil water in the morning, the white ``steam'' from the spout consists of tiny droplets condensed from water vapor---water molecules remain water molecules: a physical change. Fry an egg, and the transparent white becomes a white solid---high temperatures permanently alter protein structures: a chemical change. A cut apple turns yellowish brown after a while---molecules in its flesh react with oxygen in the air: a chemical change. Sprinkle salt into soup---the salt disappears but the soup becomes salty; its particles have merely spread out: a physical change (we will return to challenge this example later; it is not quite so simple). Charge your phone in the evening---chemical changes quietly occur inside its battery, storing electrical energy in rearranged molecules.

As you can see, there is no need to burn your textbook or visit a laboratory: your kitchen operates a chemical factory from morning to night.

\begin{shiyan}
\textbf{Materials}: two clean mineral-water bottles (about 500 milliliters each), warm water (comfortable to touch, about $35\,^{\circ}\mathrm{C}$), white sugar, a small packet of dried baking yeast (available in supermarkets), two balloons, and a marker.

\textbf{Procedure}: (1) Fill each bottle halfway with warm water, add two spoonfuls of sugar to each, and shake to mix. (2) Label the bottles: add half a spoonful of dried yeast to bottle A and shake; add nothing to bottle B. (3) Fit a balloon over each bottle opening and put the bottles somewhere warm (near a windowsill, for example, but out of strong direct sunlight). (4) Observe every 15 minutes for one to two hours, recording the changes you see.

\textbf{What to observe}: Which balloon gradually inflates? What happens to the balloon on bottle B? What gas fills the inflated balloon, and where was it hiding before? Smell bottle A after another day or two (do not drink it). How has its smell changed? Why does bottle B remain quiet---what role does it play in this setup?

\textbf{Safety}: All the ingredients are edible, but the experimental liquid has been left exposed, so do not drink it. Do not let the balloons inflate to their limit, to avoid a startling burst. Keep the water below $40\,^{\circ}\mathrm{C}$---excessive heat kills the yeast and makes the experiment ``fail.'' But that failure is itself worth thinking about: why are living things vulnerable to heat? Chapter 49, on enzymes, will answer that question.

Bottle B is the most important design feature of this experiment. Its only difference from A is ``no yeast,'' so it rules out explanations such as ``sugar water bubbles by itself'' and ``warm water bubbles by itself.'' This design is called a \textbf{control}. Chapter 2 will explain why it is essential to making experiments trustworthy.
\end{shiyan}

\section{The Chemist's Three Languages}

Chemists can describe the same event in three different languages. You need to learn to translate among them.

\textbf{The first is macroscopic language}---the language of eyes, nose, and hands: ``White sugar becomes dark and caramelized when heated, giving off a toasted aroma.'' This is the language we have been using. Its strength is direct, honest description; its weakness is that it cannot explain ``why.''

\textbf{The second is the language of microscopic particles}: ``At high temperatures, sugar molecules break down and lose water, and their atoms recombine into larger, darker molecules.'' This language explains things, but it discusses objects nobody can see---so this whole book trains your ability to picture particles in your mind.

\textbf{The third is symbolic language}, the most economical with ink and the most intimidating. Each kind of molecule can be written as a string of symbols: water is \chem{H_2O} (two hydrogen atoms and one oxygen atom), carbon dioxide is \chem{CO_2} (one carbon atom and two oxygen atoms), and glucose is \chem{C_6H_{12}O_6} (6 carbons, 12 hydrogens, and 6 oxygens). All the bustle of yeast fermentation takes just one line in symbols:

\[\chem{C_6H_{12}O_6} \rightarrow 2\chem{C_2H_5OH} + 2\chem{CO_2}\]

Read it this way: the left of the arrow is ``before the reaction,'' and the right is ``after.'' The number 2 in front means ``two molecules.'' The whole line says: yeast breaks down and reassembles one glucose molecule into two alcohol molecules and two carbon dioxide molecules.

Now do something simple enough for a small child but profoundly important: count the atoms on either side. On the left: 6 carbons, 12 hydrogens, and 6 oxygens. On the right: the two alcohol molecules contain 4 carbons, 12 hydrogens, and 2 oxygens in total; the two carbon dioxide molecules contain 2 carbons and 4 oxygens. Add them up: 6 carbons, 12 hydrogens, and 6 oxygens. \textbf{Not one extra, not one missing.} Chemical changes neither create nor destroy atoms; they merely regroup them. This line of symbols is really a ledger. Chapter 26 will teach you how to keep and balance it.

Remember: symbols are maps, not photographs. The three characters \chem{H_2O} do not tell you a water molecule's shape, size, or motion---those await Chapters 20 and 23. Symbolic language is powerful precisely because it dares to leave things out: it records only the most important information---what becomes what, and whether the atom counts match. Throughout this book, we will move repeatedly among macroscopic, microscopic, and symbolic language. When you find yourself describing a phenomenon in all three without hesitation, you are already thinking like a chemist.

\begin{qiaoliang}
Particle language says that sugar molecules ``spread out among water molecules''---but how many are there? Let us estimate. A small spoonful of glucose weighs about 5 grams. Chapter 5 will formally introduce the mole, a counting unit; for now, borrow the result that 180 grams of glucose contains about \ud{6.02\times 10^{23}}{} molecules. Then 5 grams contains approximately

\[\frac{5}{180}\times 6.02\times 10^{23}\approx 1.7\times 10^{22}\ \text{molecules}\]

This number is too large for everyday experience to offer a comparison, but try this: divide the molecules in that spoonful equally among the world's 8 billion people, and each person receives about $2\times 10^{12}$---two trillion. That is more than the total number of grains of sand on all Earth's beaches and deserts. The glass of sugar water therefore contains neither ``shadows of sugar'' nor a ``smell of sugar,'' but one hundred quintillion intact sugar molecules wandering evenly through the gaps between water molecules. Conversely, you should now see that with astronomical numbers of particles in every spoonful, chemists cannot count them individually. They must work with symbols and bulk counting methods such as the mole---the problem Chapter 5 will solve.
\end{qiaoliang}

\section{A Bowl of Caramel and a Tax Collector}

\begin{zhengju}
``Has a new substance formed?'' sounds like an easy criterion, but a hard question stands behind it: \textbf{how do you know something new has formed?} A change of color or bubbles is not enough---so what is? In the late eighteenth century, the French scientist Lavoisier answered with an unremarkable instrument: a balance.

The prevailing theory was ``phlogiston theory'': everything combustible contained a substance called phlogiston, and burning was its escape. This explained flames, but there was an awkward detail. Wood became lighter after burning (phlogiston escaped---reasonable), yet iron, tin, and lead became heavier (should they not become lighter when phlogiston escaped?). The theory's defenders had to insist that phlogiston had negative weight.

Lavoisier's approach was to lock the entire transformation ``inside the accounting office.'' He heated tin in a sealed glass vessel, first weighing the whole apparatus. Heating produced grayish white ``tin calx'' on the metal's surface. The tin had clearly changed, yet weighing the entire vessel together with the air inside showed no change whatsoever in total weight. When he opened the vessel, air rushed in with a hiss. Weighed again, the vessel was heavier, by exactly the amount by which the tin calx exceeded the original tin. There was only one conclusion: the tin's added mass came from the air.

A still more elegant experiment came around 1774. He heated mercury continuously in a closed vessel for twelve days. A red powder (then called ``mercury calx'') appeared on the silvery surface, while the volume of air in the vessel fell by about one fifth. The remaining gas extinguished candles and suffocated mice. When he collected the red powder and heated it strongly, it turned back into mercury and released a gas---whose volume exactly matched the ``missing'' fifth. This gas made glowing charcoal burn fiercely and mice lively and active. Lavoisier declared that air was not a single uniform substance, and combustion did not release phlogiston: it combined a substance with the component making up one fifth of the air, which he named oxygen.

Remember three things about this story. First, Lavoisier's ``eyes'' were his balance. He could not see molecules, but weighing let him track where matter went: color, smell, and bubbles were clues, while \textbf{precise mass accounts were evidence}. Second, he overturned phlogiston theory with a set of numbers its supporters could not explain, rather than a clever speech. The theory left the stage because its accounts could no longer accommodate ``negative-weight phlogiston,'' not because it had lost face. Third, Lavoisier himself was not entirely right: he also listed ``heat'' as an element, treating it as a substance, an error corrected only decades later. Science is not a collection of infallible geniuses; it is a long stream of checking accounts and correcting mistakes.
\end{zhengju}

\section{This Dividing Line Has Ragged Edges}

After all this discussion of ``physical or chemical change,'' it is time for an admission: the dividing line is useful, but it does not have a perfectly clean edge.

\textbf{Ragged edge one: judgment requires tools.} ``Are there new molecules?'' sounds decisive, but you cannot see or touch molecules. When sugar dissolves, what justifies saying its molecules remain unchanged? You need recovery by evaporation, mass measurements, and further instruments. The next chapter tackles how to prove that invisible things exist. Without enough evidence, the honest answer is ``we cannot tell yet,'' rather than a guess.

\textbf{Ragged edge two: some changes straddle the line.} Consider salt dissolving in water. Salt is an orderly crystal of sodium and chloride ions. In water, the crystal separates, and groups of water molecules surround each type of ion. Considerable rearrangement occurs at the particle level, yet evaporating the water returns the salt unchanged. School textbooks classify this as physical change, but you should know the classification is somewhat strained. We will reopen this account in Chapter 25, on solutions.

\textbf{Ragged edge three: some changes lie outside its jurisdiction.} A fluorescent lamp emits light when powered without producing new molecules: a physical process. In the radiation emitted by radioactive elements, however, the kinds of atoms themselves change. That goes beyond chemical change into the story of atomic nuclei in Chapter 8. Chemical change concerns ``how atoms regroup,'' not ``whether atoms themselves can change.''

\textbf{Ragged edge four: reversibility is not the criterion.} At the start, we classified our events by the intuition ``can we turn it back?'' Often useful, this intuition is no law. Charging a phone battery reverses the chemical changes of discharge; ice melting and refreezing is still physical change. Reversibility depends on energy and conditions (explained in Chapters 27 and 28), a separate question from whether a change is chemical.

Remember the proper use of this dividing line: it is a detective's magnifying glass, not a judge's gavel. Use it for a quick classification, then investigate particles and evidence more deeply whenever a case straddles the line.

\begin{wuqu}
``Bubbles, color changes, light, or heat show that a chemical change has occurred.'' This is one of the easiest traps to fall into in middle school. Refute it example by example: boiling water bubbles vigorously, but its bubbles are simply water vapor and its molecules remain unchanged---a physical change. Opening a soda bottle produces hissing bubbles as previously dissolved carbon dioxide escapes, its molecules unchanged---a physical change. An incandescent bulb is bright and hot because current heats its filament until it glows; no new substance forms---a physical change. Red ink turns a glass of water red, a dramatic color change, yet the molecules remain the same. The converse holds too: an iron railing rusts without bubbles, light, or much heat, so slowly you hardly notice, yet this is a real chemical change. Clues are calls to the police, never verdicts. To settle a case, return to the particle level and ask: are there new molecules?
\end{wuqu}

\section{Back to the Three Samples of Sugar}

Now let us keep our opening promise, view the three events through our new glasses, and check your guesses.

\textbf{Sugar water}: intact sugar molecules spread among water molecules. Unchanged molecules mean physical change. Sweetness comes from the sugar molecules themselves, so the water becomes sweet. Evaporation removes water molecules first, allowing sugar molecules to gather into crystals again, so the sugar can be recovered.

\textbf{Caramel}: at high temperatures, sugar molecules lose water, break apart, and recombine into dozens of new, dark, large and small molecules. New substances mean chemical change. Color, toasted aroma, and bitterness all come from new molecules. The building blocks have become different shapes, so the white sugar cannot be recovered.

\textbf{Sugar water with yeast}: yeast cells break down and reassemble glucose molecules into alcohol and carbon dioxide molecules, written chemically as \chem{C_6H_{12}O_6} $\rightarrow$ $2\chem{C_2H_5OH}$ $+$ $2\chem{CO_2}$. New substances mean chemical change. Carbon dioxide makes the bubbles, alcohol makes the alcoholic smell, and energy released during breakdown produces warmth. The total number of atoms remains unchanged throughout---you can check the accounts by counting as before.

And what about our opening question, which sounded almost casual: yeast is alive, so what do its activities have to do with chemistry? We can now give a serious answer. Biology studies living systems---how cells grow, reproduce, and sense their environment. Yet every action inside a yeast cell ultimately consists of molecules colliding, breaking apart, and recombining. \textbf{Life processes are not exceptions to chemistry; they are chemistry organized in an exceptionally intricate way.} This is why a book about chemistry eventually reaches cells, DNA, and evolution---the stories of Volumes II and III.

\section{Medicine Cabinets, Gardens, and an Invisible Factory}

You now hold the detective's first tool. It is time to take it out on patrol.

\textbf{The medicine cabinet.} Pick up any medicine box and read its ingredients. A common fever and pain reliever, acetaminophen, has the formula \chem{C_8H_9NO_2}. Like white sugar, it is a pure substance with a definite molecular structure; it happens to interfere with chemical reactions involved in pain signals in your body. Keep two statements firmly in mind. First, ``chemical ingredient'' does not mean ``artificial and harmful'': water, sugar, vitamin C, and the oxygen you breathe are all chemical ingredients. You yourself are a walking chemical factory made of chemicals. Second, ``natural'' does not mean ``safe'': toxins in poisonous mushrooms and sprouted potatoes are entirely natural, while lifesaving insulin is synthesized in factories. Whether a substance helps or harms depends on what it is, its dose, and who receives it, not whether its family name is ``Natural'' or ``Chemical.'' For that same reason, never adjust medication based on descriptions in a popular science book. Doses and individual differences are professional matters for doctors.

\textbf{The garden.} Soil in a flowerpot, ``nitrogen, phosphorus, and potassium'' on a fertilizer bag, reactions in sunlit leaves, and warming humus in a compost bin are all chemistry. Leaves weave carbon dioxide and water into sugar---in a sense, ``reverse engineering'' yeast fermentation. Chapter 54 will explain that precision machinery. Compost heats for the same reason as yeast sugar water: microorganisms break down molecules and release energy, doing the same trade as yeast. Every handful of fertilizer you spread supplies raw materials for invisible reactions inside plants.

\textbf{Industry.} Your breakfast bread (yeast's carbon dioxide makes its honeycomb of holes), yogurt (lactic acid bacteria turn lactose into lactic acid), soy sauce and vinegar (long microbial fermentations), monosodium glutamate, antibiotics, and a large share of modern medicines beyond monosodium glutamate all depend on choosing the right microorganisms, controlling reaction conditions, and making the right chemical changes happen. Later in Volume I, we will learn that the chemical engineer's central task is to control the direction and speed of change. What determines direction and speed is the main business of Chapters 27 through 31.

Before rushing onward, however, we must be honest about one question. Throughout this chapter we have said ``molecules,'' yet you have never seen a single one. We say sugar molecules remain in sugar water, molecules change in caramel, and the bubbles are carbon dioxide---\textbf{on what grounds?} How can we prove that something invisible and intangible exists, and that it is ``this kind'' rather than ``that kind''? This is not quibbling; it concerns the foundations of the whole scientific enterprise. In the next chapter, we become real detectives and learn how to make invisible things leave evidence of their own.

\begin{tiaozhan}
``New molecules have formed''---what is new about them? Dig one level deeper: before and after chemical change, each kind of atom retains its ``identity'' (carbon remains carbon, oxygen remains oxygen). Only two things change: the connections among atoms and the arrangement of electrons among them. The former are chemical bonds, the theme of Part III; the latter is the essence of oxidation and reduction, revealed in Chapter 33. The next question follows: why do atoms agree to recombine? The answer lies in two forces acting together---whether the new arrangement has lower energy, and whether the microscopic world's rules of probability permit it (Chapters 27 and 28). Here is another number game that might keep you awake: using only carbon, hydrogen, oxygen, and nitrogen, the theoretical number of different small molecules far exceeds the total number of atoms in the universe. You can therefore never learn chemistry by ``memorizing every substance''; you must understand its rules. Skip this passage freely if it makes no sense; it will not interrupt the main story. But if it makes you itch with curiosity, congratulations: that is exactly what this book most wants to preserve.
\end{tiaozhan}

\section*{Try It Yourself}

\begin{sikaoti}
\item Draw three particle diagrams (using small circles for molecules): (1) before and after sugar dissolves in water; (2) before and after water boils and bubbles; (3) before and after yeast sugar water bubbles. Use the diagrams to show which event changes the ``kinds of molecules.'' Add a note: circles are a representation; real molecules are not little balls.
\item Classify the following as physical or chemical changes, giving one sentence of justification for each: melting ice and snow, bread becoming moldy, perfume evaporating, an iron nail rusting, mixing flour and water into dough, and baking dough into bread.
\item A soda bottle bubbles immediately when opened; yeast sugar water bubbles gradually. What happens at the particle level in each case? Where do the molecules in each set of bubbles come from? Why can ``bubbling'' alone not prove a chemical change?
\item Count the carbon, hydrogen, and oxygen atoms on either side of this chapter's fermentation equation and check whether the ledger balances. Would it still balance if ``one glucose molecule becomes only one alcohol molecule''? What does this tell you about changing numbers in reaction equations arbitrarily?
\item A spoonful of glucose contains about $1.7\times 10^{22}$ molecules. If you could count 100 molecules per second, day and night without stopping, how many years would counting the spoonful take? (A year has about $3.2\times 10^{7}$ seconds.) Then explain why chemists need a bulk counting method such as the mole.
\item Find three things at home (nothing dangerous). For each, record the signs of change you observe, then write a sentence explaining how those clues could mislead you. How would you investigate further to confirm your interpretation?
\end{sikaoti}
